% Part: first-order-logic
% Chapter: completeness
% Section: completeness-theorem

\documentclass[../../../include/open-logic-section]{subfiles}

\begin{document}

\iftag{FOL}
      {\olfileid{fol}{com}{cth}}
      {\olfileid{pl}{com}{cth}}

\olsection{مبرهنة الاكتمال}

\begin{explain}
إذا ضممنا النتائج المتقدمة بعضها إلى بعض، تمت لنا مبرهنة الاكتمال.
\end{explain}

\begin{thm}[مبرهنة الاكتمال]
\ollabel{thm:completeness}
لتكن $\Gamma$ مجموعة من ال!!{sentence}s. إذا كانت $\Gamma$ متسقة، فهي
قابلة للإرضاء.
\end{thm}

\begin{proof}
لتكن $\Gamma$ متسقة.\iftag{FOL}{ تتيح \olref[hen]{lem:henkin} مجموعة
  متسقة مشبعة $\Gamma' \supseteq \Gamma$.}{} ثم تتيح
\olref[lin]{lem:lindenbaum} مجموعة
$\Gamma^* \supseteq \iftag{FOL}{\Gamma'}{\Gamma}$ متسقة و!!{complete}.
\iftag{FOL}{
  ومن $\Gamma' \subseteq \Gamma^*$ نعلم أن $\Gamma^*$ تضم، لكل
  !!{formula}~$!A(x)$، !!a{sentence} على الصورة
  \iftag{prvEx}
        {$\lexists[x][!A(x)] \lif !A(c)$}
        {$\lnot\lforall[x][!A(x)] \lif \lnot !A(c)$}؛
  فالمجموعة $\Gamma^*$ مشبعة كذلك. فإن كانت $\Gamma$ خالية من~$\eq$،
  أفادت}{وتفيد}
\olref[mod]{lem:truth} أن
$\iftag{FOL}{\Sat{M(\Gamma^*)}}{\pSat{v(\Gamma^*)}}{!A}$ إذا وفقط إذا
كان $!A \in \Gamma^*$. فيثبت خصوصًا لكل $!A \in
\Gamma$ أن $\iftag{FOL}{\Sat{M(\Gamma^*)}}{\pSat{v(\Gamma^*)}}{!A}$،
وهذا إرضاء لـ$\Gamma$.\iftag{FOL}{
  وأما إذا اشتملت $\Gamma$ على~$\eq$، فنستعمل \olref[ide]{lem:truth}:
  لكل ال!!{sentence}s~$!A$، يتحقق $\Sat{\equivclass{M}{\approx}}{!A}$
  إذا وفقط إذا كان $!A \in \Gamma^*$. ومن ثم
  $\Sat{\equivclass{M}{\approx}}{!A}$ لكل $!A \in \Gamma$، فيثبت
  إرضاء $\Gamma$ في هذه الحالة أيضًا.}{}
\end{proof}

\begin{cor}[مبرهنة الاكتمال، الصيغة الثانية]
\ollabel{cor:completeness}
لكل $\Gamma$ ولكل !!{sentence}~$!A$: إذا كان $\Gamma \Entails !A$،
فإن $\Gamma \Proves !A$.
\end{cor}

\begin{proof}
الرمز $\Gamma$ في كل من \olref{cor:completeness} و
\olref{thm:completeness} واقع تحت التكميم الكلي. ولرفع احتمال الخلط،
نورد \olref{thm:completeness} برمز آخر: كل مجموعة من
ال!!{sentence}s~$\Delta$ إذا اتسقت $\Delta$ كانت قابلة للإرضاء.
ومقابل ذلك بالنقيض أن $\Delta$ إذا لم تقبل الإرضاء لم تكن متسقة.
وهذه هي العبارة التي نطبقها هنا.

ليكن $\Gamma \Entails !A$. فالمجموعة $\Gamma \cup \{\lnot !A\}$
لا تقبل الإرضاء بمقتضى \olref[syn][sem]{prop:entails-unsat}. نأخذ
$\Gamma \cup \{\lnot !A\}$ مكان $\Delta$ في العبارة المتقدمة من
\olref{thm:completeness}، فيثبت عدم اتساق $\Gamma \cup \{\lnot !A\}$.
ثم بتطبيق
\iftag{FOL}{%
  \tagrefs{prfAX/{fol:axd:prv:prop:prov-incons},
    prfSC/{fol:seq:prv:prop:prov-incons},
    prfND/{fol:ntd:prv:prop:prov-incons},
    prfTab/{fol:tab:prv:prop:prov-incons}}}{%
  \tagrefs{prfAX/{pl:axd:prv:prop:prov-incons},
    prfSC/{pl:seq:prv:prop:prov-incons},
    prfND/{pl:ntd:prv:prop:prov-incons},
    prfTab/{pl:tab:prv:prop:prov-incons}}},
نحصل على $\Gamma \Proves !A$.
\end{proof}

\tagprob{FOL}
\begin{prob}
استخرج من \olref[fol][com][cth]{cor:completeness}
\olref[fol][com][cth]{thm:completeness}، ليتم بذلك إثبات تكافؤ
عبارتي مبرهنة الاكتمال.
\end{prob}
\tagendprob

\tagprob{notFOL}
\begin{prob}
استخرج من \olref[pl][com][cth]{cor:completeness}
\olref[pl][com][cth]{thm:completeness}، ليتم بذلك إثبات تكافؤ
عبارتي مبرهنة الاكتمال.
\end{prob}
\tagendprob

\tagprob{FOL}
\begin{prob}
اكتمال نسق !!a{derivation} يقتضي أن تكفي قواعده لإثبات عدم اتساق
كل مجموعة لا تقبل الإرضاء. فأي قواعد ال!!{derivation} اعتمد عليها
برهان الاكتمال؟ وأيها، إن وُجد، لم يستعمل في شيء منه؟ حقق ذلك بجدول
أو مخطط يصل قواعد ال!!{derivation} بالنتائج التي استعملتها في الطريق
إلى \olref[fol][com][cth]{thm:completeness}؛ ولا تُسقط استعمالًا
لقاعدة لمجرد أن البرهان طواه ولم يصرح به.
\end{prob}
\tagendprob

\tagprob{notFOL}
\begin{prob}
اكتمال نسق !!a{derivation} يقتضي أن تكفي قواعده لإثبات عدم اتساق
كل مجموعة لا تقبل الإرضاء. فأي قواعد ال!!{derivation} اعتمد عليها
برهان الاكتمال؟ وأيها، إن وُجد، لم يستعمل في شيء منه؟ حقق ذلك بجدول
أو مخطط يصل قواعد ال!!{derivation} بالنتائج التي استعملتها في الطريق
إلى \olref[pl][com][cth]{thm:completeness}؛ ولا تُسقط استعمالًا
لقاعدة لمجرد أن البرهان طواه ولم يصرح به.
\end{prob}
\tagendprob

\end{document}
